% =====================================================================
% 8 Cronograma
% ⚠️ O horizonte (dez/2026 a jul/2027) é SUPOSIÇÃO de trabalho, não dado do
% programa. Ajustar as colunas à data de defesa informada pelo pesquisador.
% =====================================================================

\chapter{Cronograma}
\label{cap:cronograma}

O Quadro~\ref{qua:cronograma} distribui as atividades do
Capítulo~\ref{cap:plano} por mês. \pendente{data prevista de defesa. O horizonte
abaixo, de dezembro de 2026 a julho de 2027, é suposição de trabalho e será
ajustado ao prazo do programa.} Os pedidos de acesso à informação têm prazo
legal de resposta de vinte dias, prorrogável por dez, e por isso são a primeira
atividade: o que eles devolverem condiciona o desenho de calendário.

\begin{quadro}[htb]
\centering
\caption{Cronograma de atividades}
\label{qua:cronograma}
\footnotesize
\setlength{\tabcolsep}{4pt}
\begin{tabular}{|p{7.4cm}|c|c|c|c|c|c|c|c|}
\hline
\textbf{Atividade} & \textbf{Dez} & \textbf{Jan} & \textbf{Fev} & \textbf{Mar}
  & \textbf{Abr} & \textbf{Mai} & \textbf{Jun} & \textbf{Jul} \\
 & 2026 & 2027 & 2027 & 2027 & 2027 & 2027 & 2027 & 2027 \\
\hline
Protocolo dos pedidos de acesso à informação (MAPA, ADAGRI, SEMACE)
  & X &   &   &   &   &   &   &   \\ \hline
Consulta à ação civil pública de 2011 e à Lei Municipal nº 1.511/2010
  & X &   &   &   &   &   &   &   \\ \hline
Extensão do painel até 2024 e congelamento das datas de extração
  & X & X &   &   &   &   &   &   \\ \hline
Inferência sobre o nível com 15 controles; perfil por dose; janela 2010--2014
  &   & X &   &   &   &   &   &   \\ \hline
Efeito mínimo detectável do desenho de calendário (dados anteriores à lei)
  &   & X &   &   &   &   &   &   \\ \hline
Tratamento das respostas aos pedidos de informação
  &   & X & X &   &   &   &   &   \\ \hline
Plano de análise do desenho de calendário, fechado e versionado
  &   &   & X &   &   &   &   &   \\ \hline
Estimação do desenho de calendário e robustez
  &   &   &   & X & X &   &   &   \\ \hline
Ensaio 2: leitura dos acórdãos e levantamento de custos
  &   & X & X & X &   &   &   &   \\ \hline
Redação do Ensaio 1 (versão revisada)
  &   &   &   & X & X & X &   &   \\ \hline
Redação do Ensaio 2
  &   &   &   &   & X & X &   &   \\ \hline
Introdução e conclusão gerais; revisão com o orientador
  &   &   &   &   &   & X & X &   \\ \hline
Depósito e defesa
  &   &   &   &   &   &   & X & X \\ \hline
\end{tabular}
\legend{Fonte: elaboração própria.}
\end{quadro}
